\documentclass{article}

\begin{document}
\section{Extension and cleavage}

Let $A$ and $B$ be two queries. There are six operations.
\begin{enumerate}
\item $A$ is extending $A$ independent of $B$.
\item $B$ is extending $B$ independent of $A$.
\item $\frac{A}{B}$ is extending $A$ templated by $B$.
\item $\frac{B}{A}$ is extending $B$ templated by $A$.
\item $\bar{A}$ is cleavage of $A$.
\item $\bar{B}$ is cleavage of $A$.
\end{enumerate}

The repair process is a sequence of the above six operations.
\begin{enumerate}
\item Use $\frac{A}{B}$ and $\frac{B}{A}$ as anchor.
\item $A$ and $\bar{A}$ can exchange $B$ and $bar{B}$ within each anchor (similar for $\bar{B}$).
\item $\frac{A}{B}$ and $A$ cannot be followed $\bar{A}$ until anothor $\frac{B}{A}$ comes.
\item $B$ following $\frac{A}{B}$ can be moved before $\frac{A}{B}$.
\end{enumerate}

From $\frac{A}{B}$ to the next anchor ($\frac{A}{B}$ or $\frac{B}{A}$), there can only be $A$, $B$ and $\bar{B}$. Since $\bar{B}$ cannot follow $B$, all $B$ are after $\bar{B}$. Since $A$ can exchange $B$ and $\bar{B}$, we can move all $A$ next to $\frac{A}{B}$. So the sequence within anchors are $A*\bar{B}*B*$. Since adjacent $\bar{B}$ can be merge to a single $\bar{B}$, the sequence is either $A*B*$ or $A*\bar{B}B*$. For the former, we can move all $B$ before $\frac{A}{B}$ since $B$ exchanges $A$ and $\frac{A}{B}$. This makes the anchor space as $A*$

Consider all anchor pair starts with $\frac{A}{B}$ ($\frac{B}{A}$ is symmetric).
\begin{enumerate}
\item $\frac{A}{B}X*\frac{A}{B}$. If the second $\frac{A}{B}$ is not followed by $\bar{B}$ before its next anchor, then all $B$ after it is sent to the current anchor space. However, this changes neither $A*$ nor $A*\bar{B}B*$ because for the former, we can further send all $B$ to the previous anchor space, and for the latter, these $B$ are merged to $B*$. Otherwise, the second $\frac{A}{B}$ is followed by $\bar{B}$ before its next anchor, then nothing is sent to the current anchor space.
\item $\frac{A}{B}X*\frac{B}{A}$. If $\frac{B}{A}$ is not followed by $\bar{A}$ before its next anchor, then all $A$ after it is sent to the current anchor space. These $A$ are either merged to $A*$ with or without exchanging with $\bar{B}B*$.
\end{enumerate} 

All in all, each $\frac{A}{B}$ is followed by either $A*$ or $A*\bar{B}B*$. If we do not include cleavage $\bar{B}$, then all $\frac{A}{B}$ are followed by $A*$. Before the first anchor, there cannot be cleavage. So the operator sequence before the first anchor must be $A*B*$ (or $B*A*$).

\section{Status}

\subsection{No cleavage}

The general states is $((a, i), (b, j))$.
\begin{enumerate}
\item $a$ is the current extension position of $A$.
\item $b$ is the current extension position of $B$.
\item $i$ is the pair position of $A$. It can be a position in the two references, in $B$, or $\emptyset$.
\item $j$ is similar to $i$ but for $B$. 
\end{enumerate}

We assume that $i$ and $j$ are independent. This allow us to reduce the state space to $((a, i), b)$, $(a, (b, j))$ and $(a,b)$. That is, when we extend the pair $(a, i)$, it is only matter that how long $B$ has been extended ({\it i.e.} $b$) and $j$ is always assumed to be $\emptyset$.

Here are the limitations.
\begin{enumerate}
\item $i$ and $j$ may be dependent when $A$ paired with $B$. But for simplicity, we negalect this.
\item The total states are of $O(N^3)$, which implies $O(N^3)$ iteration. 
\end{enumerate}

\section{With cleavage}

The state space is just the same as that without cleavage. However, there are two extension operations. The first the just the extention along $A/B$ (the same as that of the case without cleavage). Another is the temperary or auxilary extension. The auxilary extension of $A$ can be used by the extension (auxilary or not) by $B$ and vice versa. However, auxilary extension is only useful when it is used (maybe non-directly) by non-auxilary (permanent) extension. Because auxilary extension is allowed to be cleavaged, the with-cleavage case is hard.

\end{document}